%% ma: L11-B15-0068
%% lop: 11
%% chuong: 5
%% bai: 15
%% dang: 11.15.8
%% loai: TN4
%% muc: VD
%% dap_an: "D"
%% nguon: "tai-lieu-hoc-tap-mon-toan-lop-11-theo-chuong-trinh-giao-duc-pho-thong-2018.pdf (Lớp Toán Thầy Duy, Phạm Lê Duy, THPT Chu Văn An), Chương 3 Bài 1, Phần 1, Câu 57 (tr.333)"
%% kien_thuc: "Lớp 11 Bài 15 (giới hạn của dãy số)"
%% trang_thai: chinh-thuc
%% ghi_chu: "Đề gốc in số hạng thứ ba là 100/(10·25·30), đây là lỗi đánh máy; theo quy luật là 100/(20·25·30) nên kho đã sửa. Đáp án D khớp đáp án gốc. Đã duyệt 03/11/2026."
\begin{ex}
Tổng \[S=\dfrac{100}{10\cdot15\cdot20}+\dfrac{100}{15\cdot20\cdot25}+\dfrac{100}{20\cdot25\cdot30}+\cdots+\dfrac{100}{110\cdot115\cdot120}\] có giá trị bằng
\choice{$\dfrac{93}{1380}$}{$\dfrac{91}{13800}$}{$\dfrac9{138}$}{\True $\dfrac{91}{1380}$}
\loigiai{\sloppy Các số hạng là $\dfrac{100}{5m\cdot5(m+1)\cdot5(m+2)}=\dfrac45\cdot\dfrac1{m(m+1)(m+2)}$ với $m=2,3,\dots,22$ và $\dfrac1{m(m+1)(m+2)}=\dfrac12\left[\dfrac1{m(m+1)}-\dfrac1{(m+1)(m+2)}\right]$. Do đó \[S=\dfrac45\cdot\dfrac12\left(\dfrac1{2\cdot3}-\dfrac1{23\cdot24}\right)=\dfrac25\cdot\dfrac{91}{552}=\dfrac{91}{1380}.\]}
\end{ex}
